\chapter{Introduction}

\label{ch:introduction}

\section{Problem and Motivation}

\label{sec:motivation}

Wikidata is the largest openly licensed general-purpose knowledge graph, with more than one hundred million items.
Volunteers and bots maintain it, and it supplies structured data to Wikipedia infoboxes, search-engine knowledge panels and many other applications~\citep{vrandecic2014wikidata}.
Anyone can query it with SPARQL, but writing a correct query is hard.
Even a moderately complex question needs two kinds of knowledge: which Wikidata entities and properties are involved, and how Wikidata represents the fact in question.
The question itself usually reveals neither.
Asking ``who directed Inception'' does not tell anyone that the film is Q25188 or that the director property is P57.
Nor does it tell whether the answer sits on a direct property or behind a statement node.

Translating questions into SPARQL is therefore more than a convenience.
A language model that answers a question directly relies on what it stored in its parameters during training.
Its answer cannot be traced to a source, does not follow changes in the data, and can be wrong while sounding plausible.
An answer obtained by running a SPARQL query against Wikidata comes from a specific, inspectable source instead.
It reflects the current contents of the graph, and anyone can check it by reading the query.

Generating a query does not make the query correct, however.
A system can still produce a query that runs but answers a different question.
The difference is that such a failure can be traced to a specific step.
A query with the right identifiers but the wrong structure, for example, often returns no results, so running it already shows that something went wrong.
Execution then becomes a signal for detecting, and sometimes recovering from, failures.

The practical question is whether this works with models that are not fine-tuned for the task.
Fine-tuning a model for one knowledge graph can be effective, as the benchmark used here shows.
But it needs in-domain training data, it has to be repeated for every knowledge graph, and it is usually not possible for the frontier models that are strongest at code generation.
If frozen models can be combined into an effective pipeline instead, the approach could carry over to any knowledge graph with a suitable label index.
This thesis asks whether frozen models can do this task and, where they fail, why.

Translating a question into a Wikidata query involves two problems that can fail independently but are usually evaluated together.

The first is grounding: mapping the words in the question to Wikidata identifiers.
Inception must become Q25188, and director must become P57.
A single wrong identifier can produce wrong results or none at all.

The second is query structure: combining these identifiers into a query that expresses the question correctly.
In Wikidata this is difficult because the same fact can be stored in more than one way.
A short query pattern returns only the main value of a fact, for example who held an office.
Extra details need longer and more specific patterns: when the office was held, the unit of a quantity, or every member of a broad category.
The next chapter introduces these patterns with examples.

Because the short pattern is always valid SPARQL, a model that does not know a longer pattern is needed can use the short one and still get a query that runs, sometimes with plausible results.
A syntax check cannot catch this kind of error.
Only comparing the query's result with the expected answer can.
Because end-to-end evaluation folds grounding and structure into one score, it hides which of the two causes most errors.
That matters, since the two problems need different solutions.
Earlier error analyses have put most of the weight on entity linking.
\citet{xu2023wikisp}, for example, attribute 35.1\,\% of their system's failures to entity linking and call it ``the biggest source of errors''.
Much later work has accordingly focused on better grounding.

Separating the two stages requires a benchmark that offers more than questions and final answers.
Instruct-to-SPARQL~\citep{benamor2025instruct} provides each executable gold query together with an annotated version in which the Wikidata identifiers are replaced by human-readable labels.
Since the two versions are aligned, they yield ground truth for the intermediate linking stage directly, with no extra annotation and no external linker index.
One stage can then be held fixed while another is measured.
The benchmark authors do not propose a linking strategy of their own and leave the problem open.
This work takes up that open question.
Its main contribution is an evaluation, not a new model or linking algorithm: it combines established components into a pipeline and measures, stage by stage, where and why that pipeline fails.

\section{Research Questions}

\begin{description}
\item[RQ1]
How well do frozen frontier models translate natural-language questions into executable SPARQL over Wikidata, without fine-tuning?
\item[RQ2]
How much of the remaining error is caused by entity and predicate linking, and how much by query generation?
\item[RQ3]
Under the same circumstances, how do the three linking paradigms (fine-tuned end-to-end linking, zero-shot bi-encoder retrieval, and large-language-model reasoning over retrieved candidates) compare, and how does each stage depend on backbone scale?
\item[RQ4]
Which forms of inference-time guidance improve frozen-model query generation, and does their effect depend on whether they narrow or expand the generator's decision space?
\item[RQ5]
Can the remaining failures be detected and recovered \emph{without gold supervision}, using only signals that are available at inference time?
\end{description}

The discussion chapter answers these questions.

\section{Research Approach and
Contributions}

\label{sec:rqs}

\label{sec:contributions}

The method, stage-wise diagnostic evaluation, has four steps.
First, derive ground truth for an intermediate pipeline stage from the benchmark's own artefacts.
Second, hold all other stages fixed to isolate the stage under test.
Third, compare the conditions to attribute the resulting errors to that stage.
Fourth, test whether the identified cause responds to intervention.
The method chapter defines each step in detail.

The contributions of this thesis are:

\begin{enumerate}
\def\labelenumi{\arabic{enumi}.}
\item
\textbf{Stage-wise diagnostic evaluation}, defined as a protocol and applied end to end (derive, isolate, compare, intervene).
It includes mention-level linking ground truth derived from the benchmark's aligned queries, without new annotation or an external linker index.
\item
\textbf{Evidence-based attribution of the binding constraint to query generation (concept selection and structural composition) rather than disambiguation}.
Four linking approaches are compared on the same rows and ground truth, with the same candidate pools where a linker uses them, on two data splits, and the attribution is supported by a failure taxonomy, an analysis of Wikidata-specific constructs across the full corpus, and agreement across independent models.
The same experiments show that backbone dependence differs between linking and generation.
\item
\textbf{Controlled interventions and the execution-guided cascade}: four prompt-based interventions are tested (generic idiom guidance, schema evidence, selective application of guidance and a closed vocabulary).
They rule out several approaches and lead to a design principle, \emph{ground by constraint, not by enumeration}.
The thesis then introduces an execution-guided cascade that recovers some of the remaining errors from inference-time signals alone, without additional training, a reference query or new generation.
Together, these experiments yield a configured pipeline in which every backbone is chosen on measured performance.
\end{enumerate}

The thesis also makes a methodological contribution.
It identifies four measurement artefacts that can distort KGQA evaluation, such as a metric that rewards queries for returning nothing and silent API throttling that looks like a retrieval problem, and it shows how each can be detected.
In three of the four cases the artefact produced a more striking conclusion than the true result, which made it easy to accept.
They motivate the validation practices used throughout this work.
